\section{Framework-Defined Infrastructure: la aplicación como lenguaje de alto nivel}

\term{Framework-Defined Infrastructure} (FDI) trata las convenciones del framework como un lenguaje de alto nivel. El framework sabe qué routes son estáticas, cuáles necesitan funciones, qué datos pueden almacenarse en caché y qué solicitudes necesitan middleware; la plataforma interpreta esta semántica en build time y genera la configuración de la infraestructura subyacente. A diferencia de Infrastructure as Code, el desarrollador no declara directamente cada recurso de la nube, sino la aplicación, y la infraestructura se convierte en la salida de la compilación.

\subsection{Cloud compiler e intermediate representation}

Sea $A$ el código fuente de la aplicación, $S_F$ el extractor de semántica del framework y $C_P$ el compilador de la plataforma; entonces el plan de despliegue puede abstraerse como
\[
I=C_P\bigl(S_F(A)\bigr),
\]
donde $S_F(A)$ no es el código fuente original, sino una \term{intermediate representation} (IR, representación intermedia): route graph, static assets, function boundaries, cache policy, runtime hints y metadatos del entorno. La IR evita que la plataforma tenga que entender la implementación interna de cada framework; solo necesita un contrato de salida estable.

\lecturefigure{03-cloud-compiler.png}{La framework semantics se transforma, mediante la IR de la plataforma, en un plan de despliegue.}{Subtítulos locales 04:40--06:00; documentación oficial de Vercel sobre Framework-Defined Infrastructure.}

\begin{knowledgebox}{Lectura de la figura: por qué el framework aporta más información que una IaC genérica}
La IaC sabe «crear una function», pero no necesariamente sabe a qué route corresponde, si puede volverse estática, cuándo hacer revalidate o si requiere edge middleware. El framework conoce la semántica de la aplicación y la plataforma conoce las restricciones de ejecución; la ventaja de FDI proviene de combinar ambas en build time, no de escribir unas cuantas líneas menos de YAML.
\end{knowledgebox}

\subsection{Build Output API: la adaptación de frameworks como artifact contract}

La Build Output API de Vercel define el resultado de la compilación como una especificación del sistema de archivos: los recursos estáticos, las functions, el routing y la configuración se escriben como artifacts desplegables. Así, Next.js, SvelteKit, Nuxt u otros frameworks pueden generar una IR común, y la plataforma solo necesita consumir una salida estandarizada. También permite usar un prebuilt deployment después de hacer el build en local, con lo que el build queda desacoplado del deploy.

\lecturefigure{04-build-output-graph.png}{Build Output descompone una aplicación en un artifact graph desplegable.}{Documentación oficial de la Build Output API de Vercel; metáfora del cloud compiler presentada en clase.}

\begin{importantbox}{Cuatro beneficios de sistema que aporta la abstracción del compilador}
\begin{enumerate}
  \item \textbf{Reproducible}: las mismas entradas y versiones generan salidas que pueden auditarse;
  \item \textbf{Portable}: los frameworks se integran a la plataforma mediante el artifact contract;
  \item \textbf{Optimizable}: la plataforma puede elegir cache, runtime y region para cada route;
  \item \textbf{Reversible}: los objetos de despliegue son inmutables y el puntero de tráfico puede regresar a una versión anterior.
\end{enumerate}
\end{importantbox}

\begin{warningbox}{El compilador no puede inferir toda la intención de negocio}
La plataforma puede identificar routes y cache declarations, pero no sabe si los datos de inventario pueden estar desactualizados, si cierta API requiere consistencia fuerte ni si una falla afecta los ingresos. FDI necesita guard rails y overrides; si un valor predeterminado incorrecto se propaga automáticamente a todo el mundo, la automatización solo provoca incidentes más rápido.
\end{warningbox}

\subsection{Resumen de la sección}

La esencia de FDI es generar una IR de plataforma a partir de la semántica del framework y luego asignarla a la infraestructura. La Build Output API hace explícita esta frontera de compilación, de modo que el ecosistema de frameworks y el runtime de la plataforma pueden evolucionar de forma independiente.
